\section{Experimento}\label{ref:experimentos}

\subsection{Materiais} 

Para que o experimento fosse realizado, os materiais presentes nas listas abaixo foram utilizados.

\subsubsection{Quadripolo A}

\begin{itemize}
	\item 2 resistores de $10k\Omega$
	\item 1 resistores de $4.7k\Omega$
	\item 2 multímetros Digital GC-9205
	\item 1 protoboard
	\item 1 fonte de tensão controlada com dois canais
\end{itemize}

\subsubsection{Quadripolo B}

\begin{itemize}
	\item 2 resistores de $10k\Omega$ 
	\item 3 resistores de $20k\Omega$ 
	\item 2 multímetros Digital GC-9205
	\item 1 protoboard
	\item 1 fonte de tensão controlada com dois canais
\end{itemize}

\subsection{Métodos}

Os componentes passivos presentes no roteiro foram montados de acordo com a metodologia apresentada a seguir

\subsubsection{Quadripolo A}

Para aferir os parâmetros de impedância do quadripolo A será feita a montagem \textit{layout} do circuito pela figura \ref{QB}. Foi medido $z_{11}$ e $z_{21}$ para cada quadripolo com o multímetro seguindo o procedimento abaixo:
• Empregando uma fonte de tensão para alterar o valor $V_1$, foi ajustado a corrente $I_1$ de cada circuito para os valores sugeridos na tabela abaixo.
\begin{table}[!htb]
\caption{Valores de Corrente de Exitação}\label{TabA} 
\label{ressimA}
\centering
\begin{tabular}{c}
	$I_1$ \\
	\hline
	$0,5mA$\\
    \hline
    $0,8mA$\\
    \hline
    $1,0mA$\\
	\hline
    $1,2mA$\\
    \hline
    
	\end{tabular}
\end{table}

• Medindo as tensões $V_1$ e $V_2$ para cada conjunto de medidas   e abrindo o circuito para que $I_2 = 0$.\newline

• Invertendo a lógica fazendo $I_1 = 0$ para medir os outros parâmetros

 
 \subsubsection{Quadripolo B}
Para aferir os parâmetros de Impedância do quadripolo B será feita a montagem \textit{layout} do circuito pela figura \ref{QB}.\newline•Foi ultilizado os mesmos valores de exitação da tabela \ref{TabA}.\newline
    • Medindo as tensões $V_1$ e $V_2$ para cada conjunto de medidas   e abrindo o circuito para que $I_2 = 0$.\newline
    • Invertendo a lógica fazendo $I_1 = 0$ para medir os outros parâmetros
